\subsection{Status evaluasi dan ruang lingkup hasil}
\label{sec:exploratory-status}

Hasil berikut berasal dari eksperimen pengembangan yang telah selesai,
bukan evaluasi konfirmatori pada himpunan uji final. Label MAL mengikuti
sumber akuisisi; kandidat HN belum memiliki verifikasi benign oleh anotator
independen. Karena itu, kami melaporkan jumlah keputusan \emph{malicious},
validitas keluaran, dan perubahan keputusan, tanpa menafsirkan alarm HN
sebagai FPR atau menghitung akurasi terhadap gold. Card KB malicious/benign
terverifikasi dan pemisahan final KB--test belum selesai. Perbandingan utama
DIFF+I+MB dengan DIFF+I+MM serta ablasi intent belum memiliki hasil
konfirmatori.

Evaluasi baseline menjalankan seluruh 973 pasangan versi yang ditetapkan pada
enam model, masing-masing dengan DIFF dan DIFF+I; hasilnya ada di
Bagian~\ref{sec:baseline-models}. Baseline tersebut tidak mencakup konfigurasi
retrieval. Sebaliknya, studi versi dependency hanya mencakup 20 target
diagnostik, sedangkan studi isi dependency dan pembanding Granite
masing-masing mencakup 15 target. Pendekatan konteks dependency dan Granite
belum dievaluasi penuh pada 973 pair. Hasil subset di
Bagian~\ref{sec:dependency-diagnosis}--\ref{sec:granite-comparison} tidak
menggantikan evaluasi penuh tersebut, dan penyelesaian inferensi pada subset
tidak berarti seluruh eksperimen penelitian selesai.

\subsection{Perbandingan model pada prompt baseline}
\label{sec:baseline-models}

Tabel~\ref{tab:baseline-models} merangkum run baseline: enam model lokal,
dua konfigurasi tanpa contoh, dan kohort 973 pair Bagian~\ref{sec:run-cohort}
pada setiap kombinasi.
Seluruh 11.676 respons memakai BF16, greedy, seed 0, batas keluaran 2.048
token, konteks 32.768, dan satu permintaan pada satu waktu. Angka ini berasal
dari implementasi prompt baseline, bukan dari representasi diff, schema, dan
verifier yang dipakai sekarang. Angka pada tabel ini akan diperbarui setelah
seluruh 973 pair dijalankan ulang pada metode terkini.

\begin{table*}[t]
\caption{Enam model lokal pada 973 pair yang sama, konfigurasi tanpa contoh.
TPR adalah proporsi MAL berlabel sumber yang diberi verdict
\texttt{malicious}, bukan recall terhadap gold manusia; FPR TYP dihitung atas
100 pembaruan yang diasumsikan benign. Keluaran invalid dihitung benign.
Subset \emph{hard negative} tidak ada pada run ini.}
\label{tab:baseline-models}
\centering
\begin{tabular}{llrrrrrr}
\toprule
Model & Konfigurasi & TPR & FPR TYP & Macro kampanye & JSON valid & Median (ms) & p95 (ms) \\
\midrule
\multicolumn{8}{l}{\emph{Dalam cakupan penelitian: 3--4B}} \\
Qwen2.5-Coder-3B & DIFF   & 84,54\% & 48\% & 78,69\% & 86,84\% & 1.088 & 14.214 \\
Qwen2.5-Coder-3B & DIFF+I & 72,05\% & 49\% & 56,48\% & 75,03\% & 1.414 & 14.282 \\
Qwen3-4B         & DIFF   & 54,98\% &  6\% & 42,38\% & 71,74\% &   829 & 15.326 \\
Qwen3-4B         & DIFF+I & 63,23\% &  3\% & 54,05\% & 76,36\% &   893 & 14.114 \\
Phi-4-mini       & DIFF   & 45,93\% & 13\% & 35,87\% & 58,17\% & 1.059 & 20.168 \\
Phi-4-mini       & DIFF+I & 50,29\% & 13\% & 38,12\% & 82,94\% &   946 & 21.063 \\
\midrule
\multicolumn{8}{l}{\emph{Konteks skala, di luar cakupan: 7--8B}} \\
Qwen2.5-Coder-7B & DIFF   & 92,90\% & 65\% & 93,05\% & 94,04\% & 1.440 &  4.284 \\
Qwen2.5-Coder-7B & DIFF+I & 67,12\% & 67\% & 71,47\% & 70,81\% & 1.473 &  4.872 \\
Qwen3-8B         & DIFF   & 50,74\% &  6\% & 37,87\% & 94,96\% & 1.189 & 24.288 \\
Qwen3-8B         & DIFF+I & 48,34\% &  3\% & 33,27\% & 95,48\% & 1.520 & 24.084 \\
Llama-3.1-8B     & DIFF   & 85,45\% & 57\% & 76,44\% & 88,08\% & 1.342 & 12.025 \\
Llama-3.1-8B     & DIFF+I & 89,23\% & 69\% & 86,67\% & 89,62\% & 1.382 & 13.325 \\
\bottomrule
\end{tabular}
\end{table*}

Cakupan penelitian ini adalah model 3--4B (Tabel~\ref{tab:models}). Enam baris
pertama berada di dalam cakupan itu, dan Phi-4-mini dengan 3,836 miliar
parameter ikut masuk sebagai pembanding keluarga lain. Enam baris berikutnya
dilaporkan sebagai konteks skala, bukan sebagai klaim penelitian. Seluruh dua
belas kombinasi ditampilkan, termasuk yang menurun.

Pada rentang 3--4B tidak ada konfigurasi yang sekaligus ber-TPR tinggi dan
ber-FPR rendah. Qwen2.5-Coder-3B DIFF mencapai TPR tertinggi di kelompok itu,
84,54\%, tetapi menandai 48\% pembaruan tipikal; Qwen3-4B DIFF+I menekan alarm
tipikal ke 3\% dengan TPR 63,23\%. Pada keenam baris 3--4B, macro recall
kampanye lebih rendah daripada TPR per pair, sehingga angka per pair ikut
ditopang kampanye yang berisi banyak pair.

Arah pengaruh blok intent berbeda antar model: DIFF $\rightarrow$ DIFF+I
menaikkan TPR Qwen3-4B 8,25 poin sekaligus menurunkan FPR 3 poin, menurunkan
TPR Qwen2.5-Coder-3B 12,49 poin, dan menaikkan TPR Phi-4-mini 4,35 poin.
Kenaikan Phi tersebut terdiri dari $+$139 transisi yang melibatkan keluaran
invalid dan $-$101 pada pasangan yang kedua responsnya valid, jadi tidak dapat
dibaca sebagai perbaikan penalaran. Validitas JSON bergerak antara 58,17\% dan
86,84\% pada kelompok 3--4B dan ikut menentukan verdict, karena keluaran
invalid dihitung benign. Jika dihitung malicious, TPR/FPR Qwen2.5-Coder-3B DIFF
menjadi 96,45\%/72\%, Qwen3-4B DIFF+I 88,55\%/12\%, dan Phi-4-mini DIFF+I
67,93\%/25\%.

Tabel ini belum menjawab pertanyaan utama penelitian. Label MAL berasal dari
sumber akuisisi, bukan gold manusia; TYP diasumsikan benign; subset
\emph{hard negative} belum ada pada run ini, sehingga FPR pada pembaruan sah
yang menambahkan kemampuan sensitif tidak terukur. Interval kepercayaan belum
dihitung, konfigurasi retrieval belum dijalankan, dan tidak ada uji
signifikansi, sehingga tidak ada klaim model terbaik maupun klaim keunggulan
terhadap evaluasi lain. Latensi adalah observasi sesi dengan cache yang
berbeda, bukan tolok ukur kecepatan terkontrol.

\subsection{Diagnosis bersyarat pada versi dependency}
\label{sec:dependency-diagnosis}

Tiga bagian berikut memeriksa kecukupan konteks, yaitu faktor diagnostik di
luar enam konfigurasi Tabel~\ref{tab:configs}. Studi versi dependency meninjau
20 pembaruan yang menambahkan \texttt{easy-day-js}, dan memakai konfigurasi
DIFF+I dengan baris ber-ID serta schema JSON ketat. Kelompok ini dipilih
setelah analisis kesalahan pada dev set; seluruh target berbagi satu keluarga
dependency dan tidak dianggap sebagai 20 insiden independen.
Qwen3-4B-Instruct-2507 menerima dua snapshot dependency secara terpisah,
menghasilkan 40 inferensi baru. Instruksi kondisional pada kedua snapshot
identik. Kontrol tanpa konteks diambil dari 20 respons historis pada
konfigurasi yang sama; terhadap kontrol tersebut, penambahan instruksi
kondisional dan isi konteks tidak terpisahkan.

\begin{table}[t]
\caption{Keputusan Qwen3-4B pada 20 target studi versi dependency. Angka
merupakan jumlah verdict, bukan estimasi recall. Dua baris snapshot adalah
inferensi baru.}
\label{tab:dependency-versions}
\centering
\begin{tabular}{lrrr}
\toprule
Kondisi & Malicious & Benign & Invalid \\
\midrule
Tanpa konteks, historis & 0 & 20 & 0 \\
Snapshot 1.11.21 & 0 & 20 & 0 \\
Snapshot 1.11.22 & 20 & 0 & 0 \\
\bottomrule
\end{tabular}
\end{table}

Seluruh target berubah dari benign menjadi malicious ketika snapshot
1.11.21 diganti dengan 1.11.22 (Tabel~\ref{tab:dependency-versions}).
Versi 1.11.22 memiliki postinstall yang menjalankan \texttt{setup.cjs};
pembacaan statis file tersebut menunjukkan pengambilan konten jaringan,
penulisan file, pemanggilan proses anak, dan penonaktifan verifikasi TLS.
Sebanyak 18 dari 20 respons positif mengutip ID konteks dependency,
tetapi hanya 14 mengutip isi \texttt{setup.cjs}. ID kutipan yang sah belum
membuktikan bahwa seluruh alasan model didukung secara semantik.

Kedua versi memenuhi rentang \texttt{\textasciicircum1.11.21} dan
mendahului penerbitan target menurut metadata arsip. Namun, versi yang
benar-benar terpasang tidak diketahui karena tidak tersedia lockfile.
Hasil 20/20 malicious karenanya merupakan keputusan bersyarat pada snapshot,
bukan recall 100\%. Tidak ada keluaran terpotong atau ID bukti tak dikenal.
Total penggunaan 40 inferensi adalah 370.164 token prompt dan output.

\subsection{Kontrol konteks pada dependency lain}
\label{sec:dependency-generalization}

Studi isi dependency memakai 15 target dengan 15 nama dependency berbeda dan
mengecualikan seluruh target \texttt{easy-day-js}. Kohort terdiri dari 14
kandidat HN dan satu MAL menurut sumber, yaitu \texttt{zuper-cli@1.0.1}. Dari
67 kombinasi target--versi lokal yang memenuhi semver, 44 dikeluarkan karena
dependency diterbitkan setelah target. Dua kombinasi self-dependency, satu
versi lokal lebih lama, dan lima kombinasi yang melampaui batas nama/target
juga dikeluarkan. Versi dipilih dari arsip lokal yang tersedia dan mendahului
target; prosedur ini tidak menetapkan versi instalasi aktual atau versi
terbaru dari seluruh registry. Enam target pernah terpapar pada eksperimen
pengembangan sebelumnya; seluruh kohort tetap merupakan data pengembangan.

Kondisi tanpa isi dan kondisi dengan isi memakai instruksi kondisional serta
nama dan versi dependency yang sama. Kondisi dengan isi menambahkan isi
snapshot; kondisi tanpa isi menyatakan bahwa isi tidak tersedia. Qwen3-4B
menghasilkan 30 respons valid: seluruh 14 kandidat HN dan satu MAL dinilai
benign pada kedua kondisi. Tidak ada perubahan keputusan. Hanya dua respons
pada kondisi dengan isi yang mengutip isi dependency. Total token meningkat
dari 67.749 menjadi 128.286 (1,89 kali), tanpa perubahan verdict.

Semua kandidat HN terpilih hanya memiliki sinyal perubahan dependency;
tidak ada snapshot terpilih yang memiliki install hook. Dengan demikian,
kohort ini belum mewakili pembaruan sah dengan penambahan perilaku runtime
sensitif. Nol alarm tidak menunjukkan FPR 0\%, dan satu MAL yang tetap
benign tidak mendukung klaim kenaikan deteksi. Lima belas nama dependency
juga tidak membuktikan 15 keluarga serangan independen.

\subsection{Pembanding keluarga model: Granite 3.3 8B}
\label{sec:granite-comparison}

Pembanding Granite menjalankan IBM Granite 3.3 8B Instruct pada seluruh 15
target dan dua kondisi di atas, menghasilkan 30 inferensi baru. Teks masukan,
snapshot, schema JSON, parser, dan urutan permintaan dipertahankan.
Pembanding Qwen menggunakan 30 respons historisnya. Kedua model memakai
BF16, temperature 0, top-p 1, seed 0, batas output 2.048 token, konteks
32.768 token, dan satu permintaan pada satu waktu. Granite menggunakan
template native dengan \texttt{thinking=false} dan token akhir native;
prompt maksimum 24.467 token, tanpa pemotongan tambahan. Bobot model dan
input dipin menggunakan SHA-256 sebelum inferensi. Perbedaan keluarga,
ukuran, tokenizer, dan template berarti perbandingan ini tidak mengisolasi
efek jumlah parameter.

\begin{table*}[t]
\caption{Perbandingan pada target dan kondisi identik: Qwen3-4B (historis)
dan Granite 3.3 8B (baru). Sel menunjukkan malicious/benign/invalid.
HN belum terverifikasi benign; MAL mengikuti label sumber.}
\label{tab:granite-comparison}
\centering
\begin{tabular}{llrrr}
\toprule
Kohort & Kondisi & Target & Qwen3-4B & Granite 3.3 8B \\
\midrule
Kandidat HN & Tanpa isi & 14 & 0/14/0 & 0/14/0 \\
Kandidat HN & Dengan isi & 14 & 0/14/0 & 0/14/0 \\
MAL label sumber & Tanpa isi & 1 & 0/1/0 & 0/1/0 \\
MAL label sumber & Dengan isi & 1 & 0/1/0 & 0/1/0 \\
\bottomrule
\end{tabular}
\end{table*}

Tidak ada satu pun verdict yang berbeda antar model
(Tabel~\ref{tab:granite-comparison}). Seluruh respons Granite memenuhi
schema, tanpa keluaran terpotong, ID bukti tak dikenal, atau error transport.
Satu target MAL tetap benign, sehingga pergantian model belum memperbaiki
deteksi pada kohort ini. Dua respons pada kondisi dengan isi mengutip isi
dependency, pada \texttt{get-symbol-description@1.0.2} dan
\texttt{@tolicodes/picklejs@0.75.0}; dukungan semantik alasan belum
diverifikasi manusia.

Granite memakai 76.964 token tanpa isi dan 146.418 token dengan isi (total
223.382), dibandingkan total 196.035 token Qwen. Token antar model mengikuti
tokenizer masing-masing dan tidak disamakan dengan biaya komputasi. Median
latensi pada 14 kandidat HN adalah 1.676 ms dan 2.422 ms untuk dua kondisi
Granite; angka historis Qwen adalah 872 ms dan 1.195 ms. Latensi merupakan
observasi sesi dengan prefix caching, tanpa warmup generasi, bukan benchmark
kecepatan terkontrol. Audit pasca-inferensi mencocokkan payload, jumlah token,
parsing, kutipan, checksum sumber, artefak kontrol, dan berkas model. Hasil
ini mendukung kesamaan keputusan pada kohort yang diuji, tanpa memberikan
peringkat umum model atau bukti penurunan FPR.
